\section{Certified scalar splitting of hidden summands}
\label{sec:fitting}

The interval condition is deliberately restrictive. A visible string of
equal-matching objects may have additional maps to other objects, and those
attachments cannot be dropped. On the initial named diagram sample, the
whole-component interval test found no nontrivial eligible component. We
therefore add an exact basis-change procedure that can expose a direct sum
inside a connected differential graph. Its correctness is unconditional;
its bounded search is not complete.

\subsection{A finite linear endomorphism problem}
Let $C$ be one connected component with $B$ basis objects. Partition them by
the pair (homological degree, matching). A scalar map $Q$ may mix objects
inside each part, with coefficients in $\F$, but may not mix different parts.
Writing $Q_h$ for its degree-$h$ blocks, the chain-endomorphism equations are
\begin{equation}\label{eq:scalar-endomorphism}
 Q_{h+1}d_h=d_hQ_h.
\end{equation}
The cobordism entries have finite bit-packed bases. Equating each coefficient
in~\eqref{eq:scalar-endomorphism} produces a homogeneous binary linear system.
If the block sizes are $r_1,\ldots,r_t$, there are
$v=\sum_i r_i^2\le B^2$ unknown scalar entries. Zero and the identity are
always solutions.

This restriction to scalar maps is useful. It makes the search an ordinary
binary linear-algebra problem while still allowing nontrivial changes of
basis. It is also a limitation: a decomposition requiring more general
cobordism-valued endomorphisms can be missed.

\begin{lemma}[A cheap exact skip]\label{lem:scalar-skip}
If every (degree, matching) block contains one object and the differential
graph is connected, every scalar chain endomorphism is a scalar multiple of
the identity.
\end{lemma}
\begin{proof}
Each object has one scalar $q_v\in\F$. A nonzero entry $f:v\to w$ in
the differential gives $(q_w-q_v)f=0$. Since $f\ne0$ and the scalar field has
two elements, $q_w=q_v$. Connectivity propagates this equality to all objects.
\end{proof}
This lemma avoids solving a linear system in a common case. It does not
certify indecomposability with respect to every allowed categorical map.

\subsection{From an endomorphism to an actual split}
For a computed solution $Q$, choose a power $e$ at least the largest scalar
block dimension. The implementation uses the next power of two and repeated
squaring. In each block, form
\[
 U=\ker Q^e,\qquad V=\im Q^e.
\]
These are ordinary binary vector spaces. Assemble their bases separately in
each (degree, matching) block.

\begin{theorem}[Certified scalar Fitting split]\label{thm:fitting}
The spaces $U$ and $V$ form complementary subcomplexes of $C$. If their total
dimensions are both nonzero, changing to these bases gives a nontrivial
direct-sum decomposition of $C$ by an exact chain isomorphism in the
implemented cobordism category.
\end{theorem}
\begin{proof}
For an endomorphism of an $r$-dimensional vector space, the increasing sequence
of kernels and the decreasing sequence of images stabilize by exponent $r$.
At a stabilizing exponent, $\ker Q^e\cap\im Q^e=0$: if $z=Q^ey$ and
$Q^ez=0$, then $y\in\ker Q^{2e}=\ker Q^e$, so $z=0$. Rank-nullity gives
their direct sum equal to the original space.

Equation~\eqref{eq:scalar-endomorphism} also holds for $Q^e$. Thus the
differential maps a kernel vector to a kernel vector and an image vector to
an image vector, across every pair of matching blocks. Scalar changes of
basis inside one matching are invertible matrices of identity cobordisms,
so they are genuine isomorphisms of the additive category. Conjugating the
whole differential therefore exposes the two subcomplexes. No map incident
with the component is omitted: the component was a complete direct summand
before the operation.
\end{proof}

The original individual candidate $Q$ need not be invertible, and is not
used as the change of basis. The change-of-basis matrices $P_h$ have columns
given by the kernel basis followed by the image basis. The transformed maps
are
\begin{equation}\label{eq:conjugation}
 d'_h=P_{h+1}^{-1}d_hP_h.
\end{equation}
Every coefficient of every map crossing between the two selected sides must
vanish. This is checked explicitly. Confusing $Q$ with $P$, transforming only
one end of a differential, or ignoring an external attachment would each
make the replacement unsound.

\subsection{What is searched and what is verified}
The default pass considers components with at most 48 objects and at most
1,024 scalar variables. It inspects up to 64 nullspace-basis vectors and
16 deterministically generated linear combinations, with fixed seed 7243.
It accepts at most 16 splits in one stage. These are engineering limits, not
theorems that larger components are indecomposable. An unexamined component
is retained in full.

For every accepted split, the implementation independently checks the
coefficientwise commutator, the inverse of each basis matrix, and every
cross-block coefficient after~\eqref{eq:conjugation}. Optional witnesses
record the original matching and degree labels, the original and transformed
differentials, $Q$, $P$, the kernel dimensions, and the exponent used.
Independent tests reconstruct the conjugation using separate dense binary
matrix operations and reject deliberately corrupted witnesses.

Positive and negative search results can be cached. A key records complete
scalar-quiver data: matching names are consistently replaced by first-occurrence
indices, a uniform degree shift is removed, and every differential coefficient
and target index is retained. This preserves the linear equation system.
A cached positive result is still checked by recomputing commutation and the
transformed cross terms. Cache equality is equality of full data; a hash
collision alone cannot identify two complexes.

After a split, graph components are recomputed, and the interval and literal
sharing passes act on the exposed summands. Multiplicity ancestry is carried
through every change: a summand arising from one weighted parent inherits
that parent's weight. Basis transformations are never allowed to mix
independent weighted parents without an explicit expansion of those weights.

\subsection{Complexity and the present practical limit}
With $R$ an upper bound on the number of coefficient coordinates of a
morphism, the full scalar equations admit a conservative polynomial bound
in $B$ and $R$; ordinary dense binary elimination, for example, gives a bound
of order $RB^6$ elementary field operations. The bounded implementation is
much smaller on the tested components, but this overhead is still material.
The checkpoint theorem remains valid with this pass: scalar splitting does
not increase object count, and its cost is polynomial in the stored complex
and coefficient lengths.

A successful split can reduce later differential entries and compositions
without reducing the number of stored objects immediately. It may or may not
expose repeated templates or common-$\theta$ intervals. Section~\ref{sec:experiments}
reports all of these quantities separately. In particular, finding hidden
summands in actual knot scans does not by itself demonstrate that the
length-two quotient has been exercised on those scans.
